\pagebreak

\section{Capture a data stream - capture.3pl}\label{chap:capture}

    \index{library!capture.3pl}
    This library contains a single function. That function captures a
    data stream when triggered.
    
    This library must be incorporated by using a {\bf source} statement, i.e.
    \begin{lstlisting}[gobble=4, language=threepl]
       cource "capture.3pl"
    \end{lstlisting}



    \subsection{library modules}
    
    
        \subsubsection{capture - capture a data sequence}\label{sec:lmcapture}

            {\bf Library: capture.3pl}


            {\bf Prototype:} \\
            \vsp
            \begin{lstlisting}[gobble=12, language=threepl]
               queue(out),
               static({running, triggered}, "log")
               <-
               value(in),
               int(length),
               value(trigger_position, "uint:20", 0),
               value({start, trigger, readout}, "empty")
               ) {}
            \end{lstlisting}

            {\bf Output parameters:} \\
            {\bf out} is the output queue variable\\
            {\bf running} is true when data storage is running\\
            {\bf triggered} is true when a trigger has occured\\

            {\bf Input parameters:} \\
            {\bf in} is the data stream input.\\
            {\bf length} is the length of captured data.\\
            {\bf trigger\_position} controls the length of data captured
            after the trigger.\\
            {\bf start} is asserted to start data storing.\\
            {\bf trigger} is asserted to halt data storage after
            some number of samples.\\
            {\bf readout} is asserted to trigger stored data readout.\\

            {\bf description:} \\
            The {\bf capture()} module acts like a logic analyser. Following
            assertion of {\bf start} input data {\bf in} is stored continuously,
            the data buffer wrapping around. When {\bf trigger} is asserted data
            will continue to be stored such that the data value that occured
            coincident with the trigger is stored at the point in the stored
            data given by {\bf trigger\_position} (buffer wraparound is taken
            care of).

            The length of captured data is given by immediate input parameter
            length. This must be a power of 2 and may have values -\\
            256 to 4096 for Virtex and Virtex-E\\
            512 to 16384 for Virtex2 and Virtex2 Pro\\
            1024 to 32768 for Virtex5.

            The input data may be any type.

            Input parameter {\bf start} starts data recording. If the argument
            does not contain queue reads it must be type "log" in which case the
            transition from false to true starts recording (it may stay true for
            some time as only the transition is recognised). If the argument
            does contain queue reads then it may be any type, including "null",
            and arrival of a datum starts recording (the type and value is not
            significant). If this parameter is not supplied with an argument
            then data recording will start following global logic start and will
            continue following readout. {\bf start} is ignored if already
            running or during readout.

            Input parameter {\bf trigger} finalises the data capture, data
            recording running on for some time if necessary to position the
            captured data around the trigger position. If the argument does not
            contain queue reads it must be type "log" in which case the
            transition from false to true finalises the data capture (it may
            stay true for some time as only the transition is recognised). If
            the argument does contain queue reads then it may be any type,
            including "null", and arrival of a datum finalises the data capture
            (the type and value is not significant). If the trigger position is
            not zero then data recording must have started an appropriate number
            of clock cycles prior to the trigger to have stored valid data. The
            trigger is ignored if data recording is not running, the trigger has
            already occured or if readout is in progress.

            Input parameter {\bf readout} starts a readout of the stored data to
            the output queue. If the argument does not contain queue reads it must
            be type "log" in which case the transition from false to true starts
            the readout (it may stay true for some time as only the transition
            is recognised). If the argument does contain queue reads then it may
            be any type, including "null", and arrival of a datum starts the
            readout (the type and value is not significant). If the output queue
            blocks, readout will pause until write availability is restored. If
            this parameter is not supplied with an argument then readout will
            occur automatically following a trigger and subsequent completion of
            data capture. {\bf readout} is ignored if data recording is running.

            The input parameter {\bf trigger\_position} may have either an
            immediate or target argument. The trigger position is recognised
            each time a trigger occurs. The trigger position is the index into
            the output data block of the clock cycle in which the trigger
            occured.

            The output queue {\bf out} should be the same type as {\bf in} (it
            can differ if it can nevertheless be assigned from the type of {\bf
            in}, e.g. it may be an integer type of different width).

            Output parameter {\bf running} is true when data reecording is
            running, including after the trigger.

            Output parameter {\bf triggered} is true following the trigger
            input. It is false on completion of readout or when data recording
            is restarted if no readout is performed.
 
 
 


    \subsection{library procedures}
    
        None.




    \subsection{library functions}
    
        None.


    \subsection{library classes}
    
        None.
